% Generated by scripts/assemble.py; corrections belong in source/corrections.json.
\examhead{81}{哈尔滨工业大学 2024 年数学分析试卷}{79}
% p080-b17; source PDF page 80
\par\Needspace{5\baselineskip}
\blockanchor{p080-b17}
\majorquestion{一}下面命题正确请说明理由, 错误请举出反例.\par
% p080-b18; source PDF page 80
\blockanchor{p080-b18}
\question{1}正数列 $\{a_{n}\}, \{b_{n}\}$ 发散, 则 $\{a_{n} + b_{n}\}$ 发散.\par
% p080-b19; source PDF page 80
\blockanchor{p080-b19}
\question{2}$f(x)$ 在 $(-1, 1)$ 上可微, $f'(0) > 0$, 则存在 $\varepsilon > 0$ 满足 $f(x)$ 在邻域 $O(0, \varepsilon) \subset (-1, 1)$ 上单调.\par
% p080-b20; source PDF page 80
\blockanchor{p080-b20}
\question{3}若 $\int_{a}^{b}(f(x))^{2024}\mathrm{d}x = 0$ ，则 $f(x)\equiv 0,x\in [a,b]$\par
% p080-b21; source PDF page 80
\blockanchor{p080-b21}
\question{4}设非负项级数 $\sum_{n=1}^{\infty} a_n, \sum_{n=1}^{\infty} b_n$ 收敛, 则 $\sum_{n=1}^{\infty} \min\{a_n, b_n\}$ 收敛.\par
% p080-b22; source PDF page 80
\blockanchor{p080-b22}
\question{5}设 $f(x, y)$ 沿任意方向趋于 $(0, 0)$ 的方向极限都存在且相等, 则重极限 $\lim_{(x, y) \to (0, 0)} f(x, y)$ 一定存在.\par
% p080-b23; source PDF page 80
\par\Needspace{5\baselineskip}
\blockanchor{p080-b23}
\majorquestion{二}设 $m, n \geq 1$, 求 $\lim_{x \to 1} \left( \frac{m}{1 - x^m} - \frac{n}{1 - x^n} \right)$.\par
% p080-b24; source PDF page 80
\par\Needspace{5\baselineskip}
\blockanchor{p080-b24}
\majorquestion{三}$f(x)$ 是定义在 $(a, b)$ 上的函数, 对于任意收敛数列 $\{x_{n}\} \subset (a, b), \{f(x_{n})\}$ 收敛. 证明:\par
% p080-b25; source PDF page 80
\blockanchor{p080-b25}
(1) $f(x)\in C(a,b).$\par
% p080-b26; source PDF page 80
\blockanchor{p080-b26}
(2) $\lim_{x\to a^{+}}f(x),\lim_{x\to b^{-}}f(x)$ 存在.\par
% p080-b27; source PDF page 80
\blockanchor{p080-b27}
(3) $f(x)$ 在 $(a, b)$ 上一致连续.\par
% p080-b28; source PDF page 80
\par\Needspace{5\baselineskip}
\blockanchor{p080-b28}
\majorquestion{四}设 $f(x)$ 在 $\mathbb{R}$ 上二阶可微, 对任意的 $x \in \mathbb{R}$, $|f(x)| \leq M_0$, $|f''(x)| \leq M_2$.\par
% p080-b29; source PDF page 80
\blockanchor{p080-b29}
(1) 写出 $\forall x \in \mathbb{R}, \forall h > 0, f(x + h), f(x - h)$ 关于 $h$ 带有Lagrange余项的Taylor公式.\par
% p081-b00; source PDF page 81
\blockanchor{p081-b00}
(2) 证明: $\forall x \in \mathbb{R}, \forall h > 0, |f(x)| \leq \frac{M_0}{h} + \frac{h}{2} M_2$.\par
% p081-b01; source PDF page 81
\blockanchor{p081-b01}
(3) 证明: $\forall x \in \mathbb{R}$, $|f(x)| \leq \sqrt{2M_0M_2}$.\par
% p081-b02; source PDF page 81
\par\Needspace{5\baselineskip}
\blockanchor{p081-b02}
\majorquestion{五}解答如下问题:\par
% p081-b03; source PDF page 81
\blockanchor{p081-b03}
(1) 证明: $\lim_{b \to 1^{-}} \int_{0}^{b} \frac{\cos x}{\sqrt{1 - x^2}} \, \mathrm{d}x$ 存在.\par
% p081-b04; source PDF page 81
\blockanchor{p081-b04}
(2) 证明: $\lim_{b \to 1^{-}} \int_{0}^{b} \frac{\cos x}{\sqrt{1 - x^2}} \, \mathrm{d}x > 1$.\par
% p081-b05; source PDF page 81
\par\Needspace{5\baselineskip}
\blockanchor{p081-b05}
\majorquestion{六}设 $\{a_{n}\}$ 是单调递增的正数列, 证明:\par
% p081-b06; source PDF page 81
\blockanchor{p081-b06}
(1) 若 $\{a_{n}\}$ 有界, 则 $\sum_{n=1}^{\infty}\left(1 - \frac{a_n}{a_{n+1}}\right)$ 收敛.\par
% p081-b07; source PDF page 81
\blockanchor{p081-b07}
(2) 若 $\{a_{n}\}$ 无界, 则 $\sum_{n=1}^{\infty}\left(1 - \frac{a_n}{a_{n+1}}\right)$ 发散.\par
% p081-b08; source PDF page 81
\par\Needspace{5\baselineskip}
\blockanchor{p081-b08}
\majorquestion{七}设  $\{f_{n}(x)\}=\{xn^{\alpha}e^{-nx}\}, n=1,2,\cdots.$\par
% p081-b09; source PDF page 81
\blockanchor{p081-b09}
(1) 当 $\alpha$ 为何值时, $\{f_n(x)\}$ 在 $[0,1]$ 上收敛.\par
% p081-b10; source PDF page 81
\blockanchor{p081-b10}
(2) 当 $\alpha$ 为何值时, $\{f_n(x)\}$ 在 $[0,1]$ 上一致收敛.\par
% p081-b11; source PDF page 81
\blockanchor{p081-b11}
(3) 讨论 $\lim_{n\to \infty}\int_0^1 f_n(x)\mathrm{d}x$ 和 $\int_0^1\lim_{n\to \infty}f_n(x)\mathrm{d}x$ 的关系.\par
% p081-b12; source PDF page 81
\par\Needspace{5\baselineskip}
\blockanchor{p081-b12}
\majorquestion{八}设 $\varphi(x, y)$ 在 $(0, 0)$ 的某一邻域上连续, $f(x, y) = |x - y| \varphi(x, y)$.\par
% p081-b13; source PDF page 81
\blockanchor{p081-b13}
(1) 当 $\varphi(x, y)$ 满足什么条件时, $f_x'(0, 0), f_y'(0, 0)$ 存在?\par
% p081-b14; source PDF page 81
\blockanchor{p081-b14}
(2) 当 $\varphi(x, y)$ 满足什么条件时, $f(x, y)$ 在 $(0, 0)$ 处可微?\par
% p081-b15; source PDF page 81
\par\Needspace{5\baselineskip}
\blockanchor{p081-b15}
\majorquestion{九}求 $\lim_{n\to \infty}\frac{1}{n^4}\iiint \limits_{r\leq n}[r]\mathrm{d}x\mathrm{d}y\mathrm{d}z,$ 其中 $r = \sqrt{x^2 + y^2 + z^2}$.\par
% p081-b16; source PDF page 81
\par\Needspace{5\baselineskip}
\blockanchor{p081-b16}
\majorquestion{十}求 $\iint_{S} x^{2} \mathrm{d}y \mathrm{d}z + y^{2} \mathrm{d}z \mathrm{d}x + z^{2} \mathrm{d}x \mathrm{d}y$, 其中 $S$ 是 $z = x^{2} + y^{2}$ 从 $z = 0$ 到 $z = 1$ 的部分, 取下侧.\par
